Deepfake detection systems raise their own ethical questions. They can mitigate harm from malicious media, but they can also be misused to probe or circumvent defenses, to justify over\hyp reaching moderation, or to support surveillance infrastructures. We therefore design \textbf{MobileDeepfake} with conservative disclosure, on\hyp device processing, and auditability in mind, and we clearly separate experimental analysis from deployment guidance.

\paragraph{Dataset Governance and Gated Access.} All datasets are used under their respective licenses and terms. CelebDF\hyp v2, FaceForensics++, DeeperForensics\,1.0, and DFDC are standard research corpora widely used for detection benchmarks \cite{rossler2019faceforensicspp,li2020celebdf,dolhansky2020dfdc,jiang2020deeperforensics,verdoliva2020forensics_survey}. Deepfake\hyp Eval\hyp 2024 is a newer in\hyp the\hyp wild benchmark created to evaluate detection models on contemporary media from social platforms \cite{deepeval2024}. Access to Deepfake\hyp Eval\hyp 2024 is gated: requesters must agree not to use the dataset to develop or improve generative systems that could harm individuals, institutions, or societies, and must provide information about their affiliation and prior work in deepfake detection. The dataset uses a CC BY\hyp SA 4.0 license and may contain a small amount of NSFW content despite such content being disallowed in the original collection pipeline. In our work we respect these constraints: we treat Deepfake\hyp Eval\hyp 2024 purely as an evaluation corpus, do not redistribute the raw media, and report only aggregate statistics and plots.

\paragraph{Privacy and On-Device Processing.} The pipeline is designed to support on\hyp device inference on mobile hardware. For realistic deployments, raw images or videos need not leave the device: Stage~1 and Stage~2 can run locally via TorchScript or ONNX, and only aggregate metrics (e.g., counts of real/fake decisions, escalation rates) need to be logged. We recommend privacy\hyp preserving telemetry and minimisation of personally identifiable information if any server\hyp side components are used (e.g., batching results for monitoring). Device logs should avoid storing raw content, detailed per\hyp sample scores, or internal embeddings unless strictly necessary and properly secured. When integrating with user\hyp facing apps, operators should provide clear disclosures about what is processed locally, what is sent to servers, and what retention policies apply.

\paragraph{Dual Use and Responsible Release.} Releasing detection models and thresholds can enable benign uses (fact\hyp checking, platform moderation, research) but also malicious ones (evasion, targeted adversarial perturbations). We therefore avoid releasing sensitive threshold configurations tied to private datasets and instead describe \emph{how} thresholds are tuned (via grid search and cost\hyp sensitive objectives) without publishing fixed values for specific platforms. Where possible, we encourage operators to combine detection with other safeguards (user education, manual review, anomaly detection on model behaviour) rather than relying on a single model output. For public APIs, rate limiting, abuse monitoring, and terms prohibiting model probing are essential to reduce the risk of automated evasion attacks.

\paragraph{Bias, Fairness, and Representativeness.} Dataset composition may underrepresent certain demographics, languages, or capture conditions. Even though Deepfake\hyp Eval\hyp 2024 is significantly more diverse than earlier corpora, its creators note that it is still skewed toward English\hyp speaking, light\hyp skinned individuals from Western countries \cite{deepeval2024}. Our multi\hyp dataset training strategy partially mitigates overfitting to any single source and cross\hyp dataset evaluation highlights generalization gaps, but we do not perform a full fairness audit. In particular, we do not stratify metrics by demographic groups, language, or content domain, and our robustness sweeps focus on generic corruptions (JPEG, noise, blur, brightness) rather than group\hyp specific harms. Future work should extend this pipeline with fairness diagnostics (e.g., group\hyp disaggregated FNR/FPR, subgroup robustness curves) and, where appropriate, incorporate domain adaptation techniques designed to reduce group disparities \cite{pei2024deepfakesurvey}.

\paragraph{Societal Impact and Use Cases.} Potentially beneficial uses of \textbf{MobileDeepfake} include: (i) assisting content moderators or fact\hyp checking organisations in triaging suspicious media; (ii) providing individuals and journalists with local tools to verify suspect videos, especially in high\hyp stakes contexts such as elections or conflict reporting; and (iii) enabling researchers to study how detection performance evolves as generative models improve. At the same time, risks include: (a) false positives leading to suppression of legitimate speech or mislabelling of authentic content; (b) over\hyp reliance on automated scores to adjudicate truth in complex socio\hyp political contexts; and (c) deployment in surveillance or censorship systems without adequate oversight. We therefore recommend that any deployment of this system be accompanied by clear governance: defined policies for handling alerts, appeal mechanisms for contested decisions, and regular audits of error patterns.

\paragraph{Limitations and Human Oversight.} We acknowledge remaining failure modes, especially under severe degradations or adversarial manipulations. Our robustness experiments cover common corruptions but not targeted adversarial attacks, and our cross\hyp dataset evaluation shows that in\hyp the\hyp wild media remain challenging. From an environmental perspective, the cascade and on\hyp device inference reduce average per\hyp sample compute relative to always running a large expert model, but large\hyp scale retraining and extensive hyperparameter sweeps still incur non\hyp trivial energy and carbon costs. Future work should more systematically account for these costs when designing detection pipelines and evaluation campaigns. We recommend layered defenses and human\hyp in\hyp the\hyp loop review in high\hyp risk deployments: detectors should flag media for review rather than silently blocking it, and high\hyp impact decisions (e.g., account bans, removal of news content) should never be made solely on the basis of a single model's output. Continuous monitoring of escalation rates, false positives, false negatives, and user feedback is essential to ensure that the system's behaviour remains aligned with its intended protective role. In this sense, we view \textbf{MobileDeepfake} as a technical component within a broader socio\hyp technical response to deepfakes rather than as a standalone solution.
